\section{Introduction}
\IEEEPARstart{A}{n} electrocardiograph that has drifted out of specification still produces a
trace. A gain error of ten percent, a low-frequency corner that has moved, or a degraded
common-mode rejection do not stop the instrument; they change the amplitudes and the ST segment
that a clinician or an algorithm will read. The required performance is laid down in
IEC~60601-2-25~\cite{iec60601225}, but conformance is verified with external equipment, at type
testing and at periodic maintenance. In between, the instrument cannot know whether its analog
front-end still complies, and the deviations that appear in use are precisely those that the user
lacks the equipment to detect~\cite{chen2025}.

What the instrument can do on its own is limited. Integrated front-ends detect a detached
electrode and provide internal test signals to verify the signal chain at power-up; these are
applied inside the device and, by the manufacturer's own account, are not intended for compliance
testing~\cite{ads1298}. They do not cover the discrete network around the amplifier (protection
and bias resistors, gain-setting resistor, coupling and anti-aliasing filters, driven-right-leg
network, reference divider), nor say whether a deviation matters.

Two research communities have worked on parts of this problem without meeting. The analog test
community checks conformance to specifications, and alternate test replaces direct specification
measurements by cheap signatures from which pass or fail is inferred, with the explicit aim of
neither accepting a bad circuit nor rejecting a good
one~\cite{cheng2010,variyam1998spec,variyam1998synth,voorakaranam2007}; it targets production
test of integrated circuits. The fault-diagnosis community asks which component has failed and
answers with classifiers trained on simulated
faults~\cite{dieste2021,dieste2024,dieste2025,yang2020,rajpal2026}; there a fault is a component
that deviates from nominal by a fixed percentage, and whether the circuit is still within
specification is not asked.

For a medical instrument both questions matter, and a third appears: whether the problem is in
the circuit at all, since the skin-electrode interface is part of the signal path and varies by
orders of magnitude between subjects and electrode types~\cite{joutsen2024}. A percentage
definition of a fault is of little use here: in the front-ends studied in this paper, more than
\SI{93}{\percent} of the circuits with one component \SI{5}{\percent} off nominal, and more than
half of those with one component \SI{50}{\percent} off, still meet every requirement of the
standard. The circuit is a further issue. Chen \emph{et al.}~\cite{chen2025} have shown that the
symmetries of the usual benchmark circuits make different faults indistinguishable, and they
redesign the training circuit to remove them. An instrument designer cannot: the circuit is
given. What can be done is to find out in advance which components are indistinguishable and
report the diagnosis at that level.

This paper studies, by simulation, an in-service self-test of a single-lead ECG front-end that
uses only what the instrument can acquire with its own converter and internal test signals. Two
circuits are considered: a realistic one built around an integrated instrumentation amplifier
(\cA) and a discrete three-op-amp amplifier that stands for the circuits of the diagnosis
literature (\cB). The contributions are as follows.
\begin{enumerate}
\item A two-level self-test that joins specification verification and fault location for a
discrete front-end network. On \cA, compliance with IEC~60601-2-25 is decided with
\SI{0.9}{\percent} of escapes at \SI{7.8}{\percent} of false rejects, where a limit test on the
same measurements gives \SI{22.5}{\percent} and \SI{20.0}{\percent}.
\item A simulation study with manufacturing tolerances of a biopotential front-end of realistic
architecture: about \num{130000} Monte Carlo cases with hard, parametric, amplifier and electrode
faults, each labeled with eleven clinical specifications.
\item The separation of electrode problems from circuit faults, with gel and measured dry
electrodes~\cite{iec60601225,joutsen2024}, and the identification of its physical limit.
\item A quantified comparison of functional and percentage severity: training on component
deviation instead of on compliance raises false rejects from \SI{1.6}{\percent} to
\SI{37}{\percent}.
\item An open dataset and code.
\end{enumerate}
In addition, the groups of components that cannot be told apart are obtained from the
sensitivities of the nominal circuit, before any fault is simulated, and they explain why
component-level diagnosis does not generalize to unseen fault magnitudes while group-level
diagnosis does.

The study is based on simulation only, for two reasons. Labeling a case requires the eleven
tests of the standard, and the campaign covers some 600 fault conditions with 200 realizations
each, most of which require a component to be altered; neither is feasible on hardware at that
scale. And which faults break compliance, which can be seen and which can be told apart has to
be known before a bench experiment on a meaningful subset can be designed. Validation on hardware
is the next step.
